\documentclass[12pt]{article}

\usepackage{amsmath, amssymb}
\usepackage{geometry}
\usepackage{graphicx}
\usepackage{enumitem}

\geometry{a4paper, margin=1in}

\title{Homework Exercise 1}
\author{Write your Student ID here}
\date{Put today’s date here}

\begin{document}

\maketitle
\tableofcontents
\newpage

% -----------------------------
\section{List Practice}

\subsection{unordered list}

\textbf{(Practice 1)} There are two software packages introduced in this module:

\begin{itemize}
    \item \LaTeX
    \item Python
\end{itemize}

% -----------------------------
\subsection{ordered list}

\textbf{(Practice 2)} When it comes to math typesetting in \LaTeX, definitely there are
tons of commands and we cannot cover every detail of them in the lab learning.
Here are some tips for you:

\begin{enumerate}
    \item Learn from examples by checking online tutorial links. Compare the source
    codes and the outputs to learn new commands.

    \item For unknown commands, you can explore the TeXstudio Editor to find
    suitable buttons to insert them, or search on Internet/AI with their keywords.
    For example, ``latex commands square root''.

    \item Many math commands will be frequently used, for example, the square
    root symbol $\sqrt{x}$. Therefore, your proficiency will be greatly increased after
    three weeks’ practice.
\end{enumerate}

% -----------------------------
\subsection{nested list}

\textbf{(Practice 3)} Here is a quick overview of \LaTeX\ topics we will cover in this module:

\begin{enumerate}
    \item Lab Session 1
    \begin{itemize}
        \item Creating different types of lists
        \item Typesetting simple math
    \end{itemize}

    \item Lab Session 2
    \begin{itemize}
        \item Creating tables
        \item Inserting images
        \item Inserting programming codes
        \item Document decoration
    \end{itemize}

    \item Lab Session 3
    \begin{itemize}
        \item Typesetting complex math
        \item Aligning equations
    \end{itemize}
\end{enumerate}

\newpage

% -----------------------------
\section{Simple Math Practice}

All sub-problems in Practice 4 and 5 are organized into ordered lists.

Find your way in figuring out how to typeset those new symbols that appear
here. This will be the core part of your self-study.

% -----------------------------
\subsection{inline math}

\textbf{(Practice 4)}

\begin{enumerate}
    \item $y = \arccos x$ is the number in $[0, \pi]$ for which $\cos y = x$.

    \item In computer science, the logic not of a Boolean/logical value $p$ is commonly
    denoted by one of the following: $\sim p$, $\neg p$ or $!p$.

    \item The ceiling and floor functions of $x$ are denoted by $\lceil x \rceil$ and $\lfloor x \rfloor$.
\end{enumerate}

% -----------------------------
\subsection{display math}

\textbf{(Practice 5)}

\begin{enumerate}
    \item 
    \[
    y = 4x^3 - 2x^2 + 5x + 7
    \]

    \item 
    \[
    (1 + x)^n \approx 1 + nx + \frac{n(n - 1)}{2}x^2
    \]

    \item 
    \[
    -1 \le \sin \theta \le 1
    \]
\end{enumerate}

% -----------------------------
\subsection{inline and display math}

\textbf{(Practice 6)} The root of the quadratic equation 
$ax^2 + bx + c = 0$ is given by the quadratic formula

\[
x = \frac{-b \pm \sqrt{b^2 - 4ac}}{2a}
\]

where $\Delta = b^2 - 4ac$ is called the discriminant. The nature of roots is described
as follows:

\begin{itemize}
    \item Two distinct real roots if $\Delta > 0$.
    \item One distinct real root if $\Delta = 0$.
    \item No real root if $\Delta < 0$.
\end{itemize}

\end{document}